\documentclass[11pt]{article}
\usepackage[margin=1in]{geometry}
\usepackage{amsmath,amssymb,amsthm}
\usepackage{hyperref}
\usepackage{xurl}
\usepackage{enumitem}

\newtheorem{theorem}{Theorem}
\newtheorem{lemma}[theorem]{Lemma}
\newtheorem{corollary}[theorem]{Corollary}
\theoremstyle{definition}
\newtheorem{definition}[theorem]{Definition}
\theoremstyle{remark}
\newtheorem{remark}[theorem]{Remark}

\let\oldus\_
\renewcommand{\_}{\oldus\allowbreak}

\title{A disproof of Conjecture 00000003848\\[2pt]
\large The affine cell-core dictionary}
\date{}

\begin{document}
\sloppy
\maketitle

\section{The conjecture}

Conjecture 00000003848 reads:

\begin{quote}
\emph{Definition:} Two-sided cells of the affine 0-Hecke monoid are its two-sided equivalence classes.
\emph{Conjecture:} Two-sided cells of type $A_n^{(1)}$ correspond one-to-one with $(n+1)$-core partitions, and
the number of left cells in the corresponding cell equals the number of distinct parts of the partition
obtained by removing all $(n+1)$-hooks from that core.
\end{quote}

We show it is false. For every affine rank $n\ge 1$ (the Lean statement also covers the formal case $n=0$,
which is harmless), the empty partition is an $(n+1)$-core with $0$ distinct parts,
while every two-sided cell contains at least one left cell. For $n=1$ we also give a non-degenerate
refutation. The affine 0-Hecke monoid of type $A_1^{(1)}$ is $\mathcal J$-trivial, so every two-sided
cell contains exactly one left cell. Every staircase $(k,k-1,\dots,1)$ is a 2-core with $k$ distinct parts.
Hence, for \emph{every} bijection between cells and 2-cores, the formula fails at infinitely many non-empty
cores. Everything below is formalised in Lean 4 with Mathlib, in a single file \texttt{Basic.lean}.

\section{Reading and definitions}

\begin{definition}[Affine 0-Hecke monoid]
Let $M=(m_{ij})$ be the Coxeter matrix of type $A_n^{(1)}$ on the nodes $\{0,\dots,n\}$ (cyclic Dynkin diagram):
$m_{ii}=1$; for $n\ge 2$, $m_{ij}=3$ if $j\equiv i\pm1 \pmod{n+1}$ and $m_{ij}=2$ otherwise; for $n=1$,
$m_{01}=\infty$. The 0-Hecke monoid $H(M)$ is the monoid presented by generators $\pi_0,\dots,\pi_n$ and
relations $\pi_i\pi_i=\pi_i$ (absorption) and, for $i\ne j$ with $m_{ij}<\infty$, the braid relation
$\pi_i\pi_j\pi_i\cdots=\pi_j\pi_i\pi_j\cdots$ with $m_{ij}$ factors on each side.
In Lean: \texttt{affineA n : CoxeterMatrix (Fin (n+1))} (with $\infty$ encoded as $0$, Mathlib's convention),
the relation \texttt{ZeroHeckeRel}, built from Mathlib's \texttt{CoxeterSystem.alternatingWord}, and
\texttt{AffineZeroHecke n := PresentedMonoid (ZeroHeckeRel (affineA n))}.
\end{definition}

\begin{definition}[Cells]
In a monoid $S$, $x\,\mathcal J\,y$ iff $SxS=SyS$, i.e.\ $x=ayb$ and $y=cxd$ for some $a,b,c,d$; and
$x\,\mathcal L\,y$ iff $Sx=Sy$. These are Green's relations; the principal two-sided ideal generated by
$a$ is $S^1aS^1$ [Wikipedia, \emph{Green's relations}]. Following the conjecture's own definition, a
\emph{two-sided cell} is a $\mathcal J$-class and a \emph{left cell} is an $\mathcal L$-class. Since
$\mathcal L\subseteq\mathcal J$, each left cell lies in a unique two-sided cell. The \emph{number of left
cells} of a two-sided cell $C$ is the cardinality, in $\mathbb N\cup\{\infty\}$, of the set of
$\mathcal L$-classes contained in $C$. In Lean: \texttt{JRel}, \texttt{LRel} (both proved to be equivalence
relations), \texttt{jClass}, \texttt{lClass}, \texttt{TwoSidedCell}, \texttt{numLeftCells} (a
\texttt{Set.encard}).
\end{definition}

\begin{definition}[Hooks and cores]
Partitions are Mathlib's \texttt{YoungDiagram}s. The hook length of the cell $(i,j)$ is arm $+$ leg $+1$:
``the hook length is $a_\lambda(s)+l_\lambda(s)+1$'' [Wikipedia, \emph{Young diagram}]
(\texttt{hookLength}). A partition is a $t$-core if none of its hook lengths is divisible by $t$
(\texttt{IsCore}). For $t\ge 1$ this is equivalent to having no hook of length exactly $t$. One direction
is trivial. For the other, use a $\beta$-set $B$ of the partition: hooks of length $\ell$ correspond to pairs
$b\in B$, $g\notin B$ with $b-g=\ell$. If $t\mid \ell$, walk along $b, b-t, b-2t,\dots, g$; the first entry lies
in $B$ and the last does not, so some consecutive pair $b'\in B$, $b'-t\notin B$ occurs, giving a hook of
length exactly $t$. Moreover, our refutation uses only the empty partition and the staircases
$(k,\dots,1)$, all of whose hook lengths are odd, so both definitions give the same $2$-cores in every case used.
The number
of distinct parts is the number of distinct positive row lengths (\texttt{distinctParts}).
\end{definition}

\paragraph{The hook-removal operation.} ``Removing all $(n+1)$-hooks'' from a partition is modelled by an
arbitrary map \texttt{strip}. Theorem~\ref{thm:main} assumes only \texttt{strip}$(\varnothing)=\varnothing$,
which holds because the empty partition has no hooks at all. Theorem~\ref{thm:one} assumes only that
\texttt{strip} fixes 2-cores, which holds because a core has no hook of the given length to remove. Every
reading of ``removing all $(n+1)$-hooks'' satisfies these assumptions. We do not need to fix a model of
rim-hook removal.

\paragraph{The claim.} \texttt{CellCoreDictionary n strip} states that there is a bijection
$e$ from the two-sided cells of \texttt{AffineZeroHecke n} to the $(n+1)$-cores with
$\#\{\text{left cells in }C\}=\#\{\text{distinct parts of }\texttt{strip}(e(C))\}$ for every cell $C$.

\section{The refutation}

\begin{lemma}\label{lem:nonzero}
In any monoid, every two-sided cell contains at least one left cell (\texttt{numLeftCells\_ne\_zero}).
\end{lemma}
\begin{proof}
If $C$ is the $\mathcal J$-class of $x$, then the $\mathcal L$-class of $x$ is contained in $C$, because
$x=ay$ and $y=cx$ imply $x=ay1$ and $y=cx1$.
\end{proof}

\begin{theorem}[\texttt{not\_cellCoreDictionary}]\label{thm:main}
For every $n$ and every \texttt{strip} with \texttt{strip}$(\varnothing)=\varnothing$, the conjecture
\texttt{CellCoreDictionary n strip} is false.
\end{theorem}
\begin{proof}
The empty partition $\varnothing$ has no cells, so it is an $(n+1)$-core. If $e$ were such a bijection,
let $C=e^{-1}(\varnothing)$. Then the number of left cells in $C$ would equal the number of distinct parts of
\texttt{strip}$(\varnothing)=\varnothing$, which is $0$. This contradicts Lemma~\ref{lem:nonzero}.
\end{proof}

This argument uses nothing about the monoid beyond Lemma~\ref{lem:nonzero}. A reviewer may regard the
empty core as degenerate, so we give a second refutation that avoids it.

\begin{lemma}[\texttt{affineA1\_jTrivial}]\label{lem:jtriv}
The affine 0-Hecke monoid of type $A_1^{(1)}$ is $\mathcal J$-trivial.
\end{lemma}
\begin{proof}
Here $H=\langle \pi_0,\pi_1\mid \pi_0^2=\pi_0,\ \pi_1^2=\pi_1\rangle$, since $m_{01}=\infty$ gives no braid
relation. Let $\pi_i$ act on words over $\{0,1\}$ by prepending $i$ unless the word already starts with
$i$ (\texttt{act}). The relations hold for these maps, giving a monoid homomorphism
$\Phi:H\to\mathrm{End}(\{0,1\}^*)$ (\texttt{actHom}). Put $\mathrm{ev}(m)=\Phi(m)(\varnothing)$, the
alternating normal form. The map sending a word to the product of its generators is a left inverse of
$\mathrm{ev}$, by induction using $\pi_i\pi_i=\pi_i$. Hence $\mathrm{ev}$ is injective
(\texttt{eval\_injective}). Let $\ell(m)$ be the length of $\mathrm{ev}(m)$; note that
$\mathrm{ev}(pq)=\Phi(p)(\mathrm{ev}(q))$. Each $\pi_i$ either fixes a word or lengthens it. So
$\ell(pq)\ge\ell(q)$, with equality only if $pq=q$ (\texttt{actHom\_length}). By induction, $\mathrm{ev}(p)$
is a prefix of $\Phi(p)(w)$ for every $w$ (\texttt{eval\_prefix}). So $\ell(pq)\ge\ell(p)$, with equality
only if $pq=p$. Now suppose $x=ayb$ and $y=cxd$. Then
$\ell(y)\le\ell(ay)\le\ell(ayb)=\ell(x)\le\ell(cx)\le\ell(cxd)=\ell(y)$. So $ay=y$, then $yb=y$, and
therefore $x=y$.
\end{proof}

\begin{corollary}[\texttt{numLeftCells\_affineA1}]
For $A_1^{(1)}$, every two-sided cell is a singleton and contains exactly one left cell.
\end{corollary}

\begin{lemma}[\texttt{staircase\_isCore}, \texttt{distinctParts\_staircase}]
The staircase $\delta_k=(k,k-1,\dots,1)$ (\texttt{staircase k}, cells $\{(i,j): i+j<k\}$) is a 2-core with
exactly $k$ distinct parts.
\end{lemma}
\begin{proof}
Row $i$ has length $k-i$ and column $j$ has length $k-j$. So the hook length at $(i,j)$ is
$2(k-i-j)-1$, which is odd. The row lengths $k,k-1,\dots,1$ are distinct.
\end{proof}

\begin{theorem}[\texttt{affineA1\_failures\_infinite}, \texttt{not\_cellCoreDictionary\_one}]\label{thm:one}
Let \texttt{strip} fix every 2-core. For every bijection $e$ from the two-sided cells of the affine 0-Hecke
monoid of type $A_1^{(1)}$ to the 2-cores, the set of 2-cores $\kappa$ with
$\#\{\text{left cells in } e^{-1}(\kappa)\}\ne\#\{\text{distinct parts of }\texttt{strip}(\kappa)\}$ is
infinite. It contains every $\delta_{k}$ with $k\ge2$. In particular \texttt{CellCoreDictionary 1 strip}
is false.
\end{theorem}
\begin{proof}
The left-hand side is always $1$, and the right-hand side at $\delta_{k+2}$ is $k+2$. The map
$k\mapsto\delta_{k+2}$ is injective, because $\delta_{k+2}$ has $k+2$ distinct parts.
\end{proof}

\begin{remark}
Lemma~\ref{lem:nonzero} holds for any notion of cells in which left cells partition two-sided cells.
So Theorem~\ref{thm:main} does not depend on reading ``cells'' as Green's classes.
\end{remark}

\section{Lean formalisation and axiom audit}

Lean 4 (\texttt{leanprover/lean4:v4.33.1}) with Mathlib \texttt{v4.33.1}; file
\texttt{lean/Conjecture3848/Basic.lean} (369 lines, only \texttt{import Mathlib}). Main statements:

{\small\begin{verbatim}
theorem not_cellCoreDictionary (n : Nat)
    (strip : YoungDiagram -> YoungDiagram) (hstrip : strip bot = bot) :
    not (CellCoreDictionary n strip)
theorem affineA1_jTrivial (x y : AffineZeroHecke 1) (h : JRel x y) :
    x = y
theorem affineA1_failures_infinite (strip : YoungDiagram -> YoungDiagram)
    (hstrip : forall mu, IsCore 2 mu -> strip mu = mu)
    (e : TwoSidedCell (AffineZeroHecke 1) ~= Core 2) :
    {k : Core 2 | numLeftCells (e.symm k).1
                  != (distinctParts (strip k.1) : ENat)}.Infinite
theorem not_cellCoreDictionary_one (strip : YoungDiagram -> YoungDiagram)
    (hstrip : forall mu, IsCore 2 mu -> strip mu = mu) :
    not (CellCoreDictionary 1 strip)
\end{verbatim}}
(Shown in ASCII; the file uses the Unicode symbols.) \texttt{lake build} succeeds. \texttt{\#print axioms} for each of the four theorems reports
\texttt{[propext, Classical.choice, Quot.sound]}. There is no \texttt{sorry} and no \texttt{native\_decide}.

\paragraph{Credit.} The Green's-relation set-up follows the accepted solution of the finite analogue
00000003843 (orionsheep, GPL-3.0), generalised here to arbitrary monoids.

\begin{thebibliography}{9}
\bibitem{green} Wikipedia, \emph{Green's relations}, \url{https://en.wikipedia.org/wiki/Green\%27s_relations}.
\bibitem{young} Wikipedia, \emph{Young diagram}, \url{https://en.wikipedia.org/wiki/Young_diagram}.
\end{thebibliography}

\end{document}
